\chapter{Introduction}\label{sec:introduction}

The Phillips curve, originally proposed by \textcite{phillips1958} and later 
formalised by \textcite{samuelson1960}, suggests an inverse relationship 
between inflation and unemployment. Despite its theoretical appeal, the 
empirical stability of this relationship has been widely debated. The curve 
appeared to break down during the stagflation of the 1970s, and more recently, 
the post-COVID surge in inflation has reignited the discussion of whether 
unemployment retains meaningful predictive power over inflation dynamics 
\parencite{stock2019}. Understanding whether labour market conditions carry 
forecasting value for inflation is not merely an academic question — it has 
direct implications for monetary policy and macroeconomic stabilisation.

This paper investigates whether the Phillips curve relationship carries forecasting 
value for Danish inflation. Denmark presents an interesting case: as a small open 
economy operating under the ERM II exchange rate mechanism,\footnote{The European 
Exchange Rate Mechanism II (ERM II) is a framework under which participating 
non-euro area EU member states maintain their exchange rates within defined 
fluctuation bands against the euro. Denmark has operated under ERM II since 1999, 
effectively pegging the krone to the euro at a central rate of 7.46038 DKK per 
euro with a narrow fluctuation band of $\pm$2.25 percent. This arrangement ties 
Danish monetary conditions closely to those set by the European Central Bank.} 
Danish monetary conditions are closely tied to those of the euro area. This raises 
the question of whether broader European macroeconomic conditions carry additional 
forecasting value for Danish inflation beyond what domestic labour market variables 
can provide. We therefore extend the analysis beyond the domestic Phillips curve to 
include EU27 inflation and unemployment as potential predictors.

% Specifically, we ask: \textit{Can we produce more accurate forecasts of Danish 
% inflation by incorporating unemployment data and the broader European 
% macroeconomic environment, compared to a purely univariate benchmark?}

Specifically, the reasearch question of this paper is:

\textit{Can forecasts incorporating unemployment data and the broader European 
macroeconomic environment produce more accurate predictions of Danish inflation 
than a purely univariate benchmark?}


To answer this question, we compare five forecasting models estimated on 
monthly Eurostat data from December 2000 to March 2024, evaluated on a 
24-month out-of-sample test period from April 2024 to March 2026. A univariate 
ARIMA model serves as the benchmark, and is also estimated with structural break 
dummies to assess whether controlling for identified level shifts improves forecast 
accuracy. Against these two ARIMA specifications we compare three dynamic 
regression models incorporating Danish unemployment and EU27 inflation as 
external predictors. Variable selection is guided by a General-to-Specific (GETS) approach to 
account for structural instability and outliers in the data 
\parencite{doornik2014}. 

% Two structural breaks in Danish 
% inflation, identified at December 2012 and August 2021, are incorporated 
% throughout the analysis.

% The results show that the univariate ARIMA benchmark outperforms every 
% multivariate specification out-of-sample. The GETS procedure eliminates all 
% Danish unemployment variables as statistically insignificant, providing direct 
% evidence against the practical forecasting value of the Phillips curve in this 
% sample. EU27 inflation is retained by GETS and is highly significant in-sample, 
% confirming the ERM II spillover channel, but it does not translate into improved 
% out-of-sample forecasts during the post-surge normalisation period. These 
% results are consistent with the broader finding in the inflation forecasting 
% literature that multivariate models frequently fail to beat simple univariate 
% benchmarks \parencite{stock_watson_2008}.

The remainder of the paper is structured as follows. Section 2 describes the 
data. Section 3 introduces the methodology. Section 4 presents and discusses 
the results. Section 5 concludes.